\begin{table}[!ht]
\caption{Resource comparison: scikit-learn HDBSCAN vs.\ GSHAC on a 100K
  subset of the GeoNames sample (\(n = 100{,}000\), haversine), both measured
  in this run with the same whole-pipeline peak-memory definition.
  Non-equivalent outputs; illustrates resource profiles only. The gap widens with
  \(n\): HDBSCAN scales super-linearly (extrapolated to hours at 2M), GSHAC
  near-linearly (Table~\ref{tab:geonames}).}
\label{tab:hdbscan}
\small
\begin{tabular}{lrrrr}
\toprule
\textbf{Method} & \textbf{Time} & \textbf{Peak mem} & \textbf{Clusters} & \textbf{Noise} \\
\midrule
HDBSCAN (mcs\,=\,2) & 229\,s & 43\,MiB & 29{,}398 & 12.7\% \\
HDBSCAN (mcs\,=\,5) & 219\,s & 33\,MiB & 4{,}667 & 35.2\% \\
\addlinespace
GSHAC (\(\hmax\!=\!10\)\,km)\(^{*}\) & 5.6\,s & 35\,MiB & --- & 0\% \\
\quad\(\to h\!=\!5\)\,km  & & & 80{,}571 & \\
\quad\(\to h\!=\!10\)\,km & & & 56{,}141 & \\
\bottomrule
\multicolumn{5}{l}{\footnotesize \(^{*}\)Both: mean over 5 repetitions. One dendrogram cut at any \(h \leq \hmax\).} \\
\multicolumn{5}{l}{\footnotesize GSHAC peak is whole-pipeline: graph phase 35\,MiB, MST-path HAC phase 24\,MiB.}
\end{tabular}
\end{table}